% -*- TeX:UK -*-
\documentclass[10pt]{beamer}
\usetheme{metropolis}
%\useinnertheme{rectangles}'
\setbeamercovered{%
still covered={\opaqueness<1->{15}},
again covered={\opaqueness<1->{40}}}

\hypersetup{colorlinks,linkcolor=black,urlcolor=brown,citecolor=brown}

\usepackage{wasysym}
\usepackage{amsmath,amssymb,amsthm}
\usepackage{unicode-math}

% We set the Lucida OTF fonts as default
\usepackage{fontspec}
\setmainfont{Lucida Bright OT}
\setsansfont{Lucida Sans OT}
\setmonofont{Lucida Console DK}[Scale=MatchLowercase]

\newfontfamily\typiconsfont{Typicons}[Scale=0.9]
\newcommand\typicons[1]{{\typiconsfont\symbol{#1}}}
\newcommand\Attention{\colorbox{yellow}{\textcolor{black}{\typicons{"E136}}}\xspace}
\newcommand\Advanced{\colorbox{black}{\textcolor{white}{\typicons{"E04E}}}\xspace}

\newfontfamily\webglyphsfont{WebHostingHub-Glyphs}[Scale=0.7]
\newcommand\webglyphs[1]{{\webglyphsfont\symbol{#1}}}
\newcommand\Discussion{\colorbox{white}{\textcolor{black}{\webglyphs{"F134}}}\xspace}
\newcommand\DExamples{\colorbox{black}{\textcolor{white}{\webglyphs{"F134} examples?}}}
\newcommand\Reading{\colorbox{black}{\textcolor{white}{\webglyphs{"F0C1}}}\xspace}

\newfontfamily\meteoconsfont{Meteocons}
\newcommand\meteocons[1]{{\meteoconsfont\symbol{#1}}}
\newcommand\meteosun{\meteocons{"0042}}
\newcommand\meteosuncloud{\meteocons{"0048}}
\newcommand\meteorain{\meteocons{"0052}}
\newcommand\meteosolidsun{\meteocons{"0031}}

\newfontfamily\KRfarmfont{KR Back On The Farm}
\newcommand\KRfarm[1]{{\KRfarmfont\symbol{#1}}}
\newcommand\farmplant{\KRfarm{"0049}}

\usepackage{polyglossia}
\setdefaultlanguage[variant = british, ordinalmonthday = false]{english}

\usepackage[style=authoryear-comp,firstinits,sortcites,maxcitenames=2,%
    mincitenames=1,maxbibnames=10,minbibnames=10,uniquename=mininit,%
    uniquelist=minyear,sortfirstinits=true]{biblatex}
\addbibresource{../references/ecophys.bib}
%\renewcommand{\bibfont}{\small}

\usepackage{abbrev}
\DeclareGraphicsExtensions{.png,.jpg,.jpeg,.eps,.pdf}

\begin{document}

\title{PBIO-141\\Sensory and Physiological Ecology of  Plants}
\subtitle{12a: Carbon Assimilation (D)}
\author{Pedro J. Aphalo}
\date{January-February 2018}
\institute[Univ.\ of Helsinki]{M.Sc.\ in Plant Biology, University of Helsinki\\[2ex] \url{http://blogs.helsinki.fi/aphalo/}}

  \begin{frame}
    \maketitle
  \end{frame}

  \begin{frame}[c]
    \begin{center}
      \begin{small}
        \copyright 2006-2018 by Pedro J. Aphalo\\
       University of Helsinki, Finland.\\
        \textcolor{blue}{\url{http://blogs.helsinki.fi/aphalo/}}\\[2ex]
      \end{small}

      \begin{footnotesize}
        Sensory and Physiological Ecology of  Plants: `12a: Carbon Assimilation (D)' by Pedro J. Aphalo is licensed under a Creative Commons Attribution-ShareAlike 4.0 International License.

        Illustrations and text quoted from copyrighted sources is excluded from this license and their use should respect the original licenses.\\[2ex]
      \end{footnotesize}

      \centering\includegraphics[width=6em]{../figures/copyright/by-sa}
    \end{center}
  \end{frame}

  \begin{frame}
    \frametitle{Outline}
    \tableofcontents
  \end{frame}

\section{Leaf age and C assimilation}

\begin{frame}{\COtwo\ in the leaf}
    \centering
    \includegraphics[height=0.70\textheight]{figures/Niinemets3_A_vs_Cc_Ci.eps}\\
    {\small \autocite[from][]{Niinemets2005b}.}
\end{frame}

\begin{frame}{A \textit{vs.}\ leaf age}
    \centering
    \includegraphics[height=0.85\textheight]{figures/Niinemets4_A_vs_Leaf_age.eps}\\
    {\small \autocite[from][]{Niinemets2005b}.}
\end{frame}

\begin{frame}{Rubisco and N \textit{vs.}\ growth irradiance}
    \centering
    \includegraphics[height=0.85\textheight]{figures/NiinemetsB4_RubiscoNMass_vs_Leaf_age.eps}\\
    {\small \autocite[from][]{Niinemets2007c}.}
\end{frame}

\begin{frame}{Mesophyll conductance and irradiance}
    \centering
    \includegraphics[height=0.70\textheight]{figures/NiinemetsB8_gm_vs_I_and_Leaf_age.eps}\\
    {\small \autocite[from][]{Niinemets2007c}.}
\end{frame}


\section{Stable isotopes of carbon}

\begin{frame}{Isotopes of C}
    \begin{description}
        \item[$^{11}$C] radioactive, half-life: 20.3 min
        \item[$^{12}$C] stable, natural abundance: 98.9\%
        \item[$^{13}$C] stable, natural abundance:  1.1\%
        \item[$^{14}$C] radioactive, half-life: 5730 years
        \item There are also other shorter lived isotopes.
    \end{description}
\end{frame}

\begin{frame}{Radioactive isotopes of C}
    \begin{itemize}
        \item Used to trace the fate of the products of photosynthesis and photosynthetic pathways by looking how long after pulse labelling with $^{14}$\COtwo the $^{14}$C appears in different metabolites.
        \item Used to trace metabolic routes.
        \item Used to trace degradation of metabolites.
        \item Earlier used to measure photosynthetic rate.
    \end{itemize}
\end{frame}

\begin{frame}{Stable isotopes of C}
    \begin{itemize}
        \item As all isotopes they differ in their atomic mass (and size)
        \item This affects diffusion rates and also affinity to enzymes
        \item The selectivity is a property of each enzyme
        \item Final effect will depend also on concentrations
        \item Isotope discrimination in C3 plants depends on $C_i$
        \item Isotope discrimination is different in C3 and C4 plants
        \item Metabolites and other compounds containing different isotopes can
        be quantified by mass spectroscopy (MS).
        \item Even location of the isotope within molecules can be usually
        detected (in fragments).
    \end{itemize}
\end{frame}

\begin{frame}{Magnitude of fractionation}
    \begin{itemize}
        \item PEP carboxylase: 2.2 \permil
        \item Rubisco: 30 \permil\ at 25\,C
        \item[]
        \item Rubisco discriminates against $^{13}$CO$_2$
    \end{itemize}
\end{frame}

\begin{frame}{Isotope composition vs.\ fractionation}
    \begin{itemize}
        \item Composition:\\
        $\delta ^{13}$C(\permil)$ = (R_\mathrm{sample}/R_\mathrm{standard} - 1) \times 1000$
        \item Fractionation:\\
        $\Delta = \delta_\mathrm{source} - \delta_\mathrm{plant}$
        \item $\delta_\mathrm{air} \approx -8$\permil; seasonal variation in 1982: -8.0 to -7.2 \permil; in 2008: -8.7 to -7.8 \permil; in addition $\delta_\mathrm{air}$ also varies within canopies.
        \item $R$ ratio between $^{13}$C and $^{12}$C
        \item $R_\mathrm{air} \approx 0.011147$ and it is decreasing because of fossil fuel use
        \item $R_\mathrm{standard} = 0.0112372$ in Pee Dee Belemnite or (PDB) calcareous marine fossil
    \end{itemize}
\end{frame}

\begin{frame}{Carbon isotope ratio, C3 and C4}
    \centering
    \includegraphics[height=0.85\textheight]{figures/Ehleringer13_4delta13CC3C4.eps}\\
    {\small \autocite[from][]{EhlOsm1991}.}
\end{frame}

\begin{frame}{C isotope ratio and $C_i$ (C3 species)}
    \centering
    \includegraphics[height=0.85\textheight]{figures/Ehleringer13_7delta13CCi.eps}\\
    {\small \autocite[from][]{EhlOsm1991}.}
\end{frame}

\begin{frame}{Carbon isotope ratio and $C_i$}
    \centering
    \includegraphics[height=0.80\textheight]{figures/Lambers1Delta13CCiC3C4.eps}\\
    {\small \autocite[from][]{LambersEtAl1998}.}
\end{frame}

\begin{frame}{Carbon isotope ratio and irradiance}
    \centering
    \includegraphics[height=0.8\textheight]{figures/Ehleringer13_3delta13CPAR.eps}\\
    {\small \autocite[from][]{EhlOsm1991}.}
\end{frame}

\begin{frame}{Revisiting the mountain}
    \centering
    \includegraphics[height=0.85\textheight]{figures/Ehleringer13_6delta13CMountain.eps}\\
    {\small \autocite[from][]{EhlOsm1991}.}
\end{frame}

\begin{frame}{Carbon isotope ratio and seasons}
    \centering
    \includegraphics[height=0.85\textheight]{figures/Ehleringer13_9delta13CMonths.eps}\\
    {\small \autocite[from][]{EhlOsm1991}.}
\end{frame}

\section{Respiration}

\begin{frame}{Respiration}
    \begin{itemize}
        \item Respiration is essentially the reverse of
        photosynthesis.
        \item In photosynthesis C from \COtwo is reduced to
        carbohydrates, releasing O$_2$, and using energy from
        sunlight.
        \item In respiration C from carbohydrates is oxidized to
        \COtwo, releasing energy, and using O$_2$ as final acceptor
        of the electrons.
        \item In respiration about 50\% of the energy stored in
        the carbohydrates is conserved as high energy bonds in ATP.
        \item This is assuming that all C in sucrose is oxidized,
        in practice many intermediate metabolites are the start of
        biosynthetic routes.
    \end{itemize}
\end{frame}

\begin{frame}{Respiration: the fuel}
    \begin{itemize}
        \item Although respiration is most frequently described
        as starting from sucrose in plants and glucose in animals,
        it is a much more flexible process than that.
        \item Products from starch hydrolysis in plastids can
        enter the respiratory pathway.
        \item Organic acids can enter at later stages, or even
        products of the degradation of lipids.
    \end{itemize}
\end{frame}

\begin{frame}{Respiration: carbon skeletons}
    \begin{itemize}
        \item Quite often the intermediate metabolites of the
        respiratory pathway are
        used as substrates (``carbon skeletons'') in other biosynthetic pathways.
        \item Different metabolic pathways do not exist in
        isolation in the cell, although compartmentalization in
        organelles provides some degree of separation.
        \item In different cell types these connections between pathways
        will differ in their importance.
    \end{itemize}
\end{frame}

\begin{frame}{Respiration: metabolism}
    \begin{itemize}
        \item Glycolysis: sugars $\rightarrow$ organic acids (+ ATP).
        \item Pentose phosphate pathway: hexose-P $\rightarrow$
        triose-P (+ \COtwo + NADPH).
        \item Citric acid cycle: organic acids $\rightarrow$
        (\COtwo + NADH + FADH$_2$).
        \item Oxidative phosphorylation: NADH, FADH$_2$, O$_2$
        $\rightarrow$ ATP.
    \end{itemize}
\end{frame}

\begin{frame}{Respiration: metabolism}
    \centering
    \includegraphics[height=0.9\textheight]{figures/Taiz11.1Respiration.eps}\\
    {\small \autocite[from][]{TaiZei2006}.}
\end{frame}

\begin{frame}{Energy metabolism}
    \centering
    \includegraphics[height=0.8\textheight]{figures/Heldt5.1MitochondrialEnergyMetabolism.eps}\\
    {\small \autocite[from][]{Heldt1997}.}
\end{frame}

\begin{frame}{Oxidative phosphorylation}
    \begin{itemize}
        \item It is a controlled oxidation through an electron
        transport chain.
        \item Electron transport is accompanied by the pumping of
        protons (H$^+$) from the mitochondrial matrix into the intermembrane
        space of the mitochondrion.
        \item The proton gradient then drives ATP synthesis on
        the matrix side of the inner membrane (1 ATP for each 3
        protons ``returning'' to the matrix).
    \end{itemize}
\end{frame}

\begin{frame}{Alternative oxidase}
    \begin{itemize}
        \item In the mitochondrial membrane there is an
        alternative path for electron transport.
        \item The alternative oxidase accepts electrons from ubiquinone
        without pumping of protons.
        \item This partly decouples electron transport from ATP
        synthesis, generating much less ATP and more heat as a
        result.
        \item Its activity is regulated by the reduction level of the ubiquinone
        pool and by pyruvate. Transcription is induced by various stress factors.
        \item Metabolic heat generation is rare in plants, but exists in
        the flowers of some species.
    \end{itemize}
\end{frame}

\begin{frame}{Heat production}
    \centering
    \includegraphics[height=0.85\textheight]{figures/Lambers10TemperatureAndRespirationLotus.eps}\\
    {\small \autocite[from][]{LambersEtAl1998}.}
\end{frame}

\begin{frame}{Respiration: acclimation}
    \centering
    \includegraphics[height=0.8\textheight]{figures/Lambers17RespirationvsTVaccinium.eps}\\
    {\small \autocite[from][]{LambersEtAl1998}.}
\end{frame}

\begin{frame}{Photosynthate fraction respired}
    \centering
    \includegraphics[width=\textwidth]{figures/Lambers20PhotosynthateFractionRespired.eps}\\
    {\small \autocite[from][]{LambersEtAl1998}.}
\end{frame}

\begin{frame}{Construction costs}
    \centering
    \includegraphics[height=0.8\textheight]{figures/Lambers21ConstructionCosts.eps}\\
    {\small \autocite[from][]{LambersEtAl1998}.}
\end{frame}

\begin{frame}{Construction costs}
    \centering
    \includegraphics[height=0.8\textheight]{figures/Lambers24ConstructionCosts.eps}\\
    {\small \autocite[from][]{LambersEtAl1998}.}
\end{frame}

\begin{frame}{Construction and maintenance}
    \centering
    \includegraphics[height=0.65\textheight]{figures/Lambers25AnnualCUse.eps}\\
    {\small \autocite[from][]{LambersEtAl1998}.}
\end{frame}

\begin{frame}{Fate of carbon}
    \centering
    \includegraphics[height=0.8\textheight]{figures/Lambers22FateOfCarbon.eps}\\
    {\small \autocite[from][]{LambersEtAl1998}.}
\end{frame}

  \section*{References}
  \begin{frame}[t,allowframebreaks]
    \frametitle{References}
    \printbibliography
  \end{frame}

\end{document}

